% AMTTP UK Patent Application — Specification
% Title: A Multi-Layer Transaction Processing System for Compliance Enforcement
%        in Decentralised Finance Networks
% Applicant: Odeyemi Olusegun Israel
% Date: March 2026

\documentclass[12pt,a4paper]{article}
\usepackage[margin=2.5cm]{geometry}
\usepackage{amsmath,amssymb}
\usepackage{graphicx}
\usepackage{enumitem}
\usepackage{hyperref}
\usepackage{titlesec}
\usepackage{fancyhdr}
\usepackage{booktabs}
\setlength{\headheight}{14.5pt}

\pagestyle{fancy}
\fancyhf{}
\rhead{Patent Application}
\lhead{AMTTP --- Odeyemi O.I.}
\rfoot{Page \thepage}

\titleformat{\section}{\large\bfseries}{\thesection.}{0.5em}{}
\titleformat{\subsection}{\normalsize\bfseries}{\thesubsection.}{0.5em}{}

\begin{document}

\begin{center}
{\LARGE\bfseries PATENT SPECIFICATION}\\[1em]
{\Large A Multi-Layer Transaction Processing System for Compliance\\
Enforcement in Decentralised Finance Networks}\\[2em]
{\large Applicant: Odeyemi Olusegun Israel}\\[0.5em]
{\large Inventor: Odeyemi Olusegun Israel}\\[0.5em]
{\large Date of Filing: 4 March 2026}\\[0.5em]
{\large Application Number: [TO BE ASSIGNED]}\\[0.3em]
{\normalsize \textit{Filed pursuant to the Patents Act 1977}}
\end{center}

\vspace{2em}

% ============================================================
\section{TECHNICAL FIELD}
% ============================================================

The present invention relates to the field of financial technology, and more 
particularly to systems and methods for enforcing anti-money laundering (AML) 
and regulatory compliance in decentralised finance transactions using 
blockchain smart contracts, machine learning risk scoring, and 
zero-knowledge cryptographic proofs.

% ============================================================
\section{BACKGROUND OF THE INVENTION}
% ============================================================

\subsection{Technical Problem}

Decentralised finance (DeFi) protocols process financial transactions on 
blockchain networks without centralised intermediaries. While this 
architecture offers advantages in transparency and accessibility, it creates 
significant regulatory compliance challenges. Current compliance approaches 
suffer from several deficiencies:

\begin{enumerate}[label=(\alph*)]
\item \textbf{Retroactive enforcement:} Existing blockchain analytics 
platforms (e.g., Chainalysis, Elliptic) operate as retrospective monitoring 
tools. They identify suspicious transactions \emph{after} settlement, by 
which time value transfer has already occurred and may be irrecoverable.

\item \textbf{Platform-level rather than protocol-level compliance:} 
Conventional KYC/AML checks are performed at the exchange or platform level. 
When a user transacts through a DeFi protocol directly, these checks are 
bypassed entirely. Compliance is bound to the platform, not to the 
transaction.

\item \textbf{Reliance on labelled training data:} Supervised fraud detection 
models require large quantities of labelled fraudulent transactions. In 
practice, fewer than 0.1\% of blockchain transactions are labelled as 
illicit, creating severe class imbalance and data scarcity problems.

\item \textbf{Privacy-compliance tension:} Compliance verification 
traditionally requires disclosure of sensitive personal and financial data. 
This conflicts with the pseudonymous nature of blockchain transactions and 
with data protection regulations such as the UK Data Protection Act 2018 and 
the EU General Data Protection Regulation.
\end{enumerate}

\subsection{Prior Art Limitations}

Prior art systems known to the applicant include:

\begin{itemize}
\item Blockchain analytics platforms that provide post-hoc transaction 
tracing and address clustering but do not enforce compliance at the point of 
transaction.

\item Exchange-level compliance systems that implement KYC/AML at 
centralised access points but cannot extend to protocol-level DeFi 
transactions.

\item Standard machine learning fraud detection systems that require 
extensive labelled datasets and do not integrate with blockchain smart 
contract enforcement mechanisms.

\item Zero-knowledge proof systems used in cryptocurrency privacy (e.g., 
Zcash) that conceal transaction details but do not implement compliance 
verification.
\end{itemize}

None of the prior art systems provides a unified architecture that combines 
pre-settlement compliance enforcement, weak supervision for label-scarce 
environments, graph-based analytics, sanctions screening, deterministic 
policy enforcement via smart contracts, and privacy-preserving verification 
via zero-knowledge proofs.

% ============================================================
\section{SUMMARY OF THE INVENTION}
% ============================================================

The present invention provides a multi-layer transaction processing system 
that enforces regulatory compliance \emph{before} on-chain settlement in 
decentralised finance networks. The system addresses the limitations of the 
prior art through the following technical contributions:

\begin{enumerate}[label=(\roman*)]
\item A four-layer architecture comprising interface, compliance 
orchestration, offline training, and infrastructure layers that together 
implement a deterministic compliance lifecycle.

\item A weak supervision machine learning pipeline comprising a composite 
teacher ensemble that generates pseudo-labels from unlabelled blockchain 
transaction data, enabling the training of a student model without requiring 
labelled fraud data.

\item A deterministic compliance decision matrix that maps continuous risk 
scores to discrete policy actions (approve, review, escrow, block) via fixed 
thresholds, ensuring reproducible and auditable decisions.

\item A zero-knowledge proof subsystem (zkNAF) comprising three proof 
circuits---sanctions non-membership, risk range, and KYC verification---that 
enable privacy-preserving compliance verification on-chain.

\item Smart contract infrastructure that enforces compliance decisions 
deterministically, including automatic escrow of high-risk transactions and 
blocking of transactions exceeding risk thresholds.

\item A state-adaptive transaction flow control module configured to compute 
network instability metrics and to select, from a predetermined set of 
enforcement profiles, one or more enforcement parameters that modulate the 
deterministic compliance decision process in response to the computed system 
state.
\end{enumerate}

% ============================================================
\section{DETAILED DESCRIPTION OF THE INVENTION}
% ============================================================

\subsection{System Architecture Overview}

The system comprises four hierarchically organised processing layers, 
described below from the outermost (user-facing) to the innermost 
(infrastructure):

\textbf{Layer I --- Interface Layer.} This layer provides the entry points 
through which external systems interact with the compliance system. It 
comprises:
\begin{itemize}
\item Client software development kits (SDKs) in TypeScript and Python that 
expose methods for submitting transactions, querying risk scores, and 
verifying compliance status.
\item A RESTful application programming interface (API) providing HTTP 
endpoints for transaction submission, risk score retrieval, compliance 
status queries, and administrative functions.
\item A web-based dashboard application providing visual interfaces for 
compliance monitoring, transaction review, and system administration.
\end{itemize}

\textbf{Layer II --- Compliance Orchestration Layer.} This layer implements 
the core compliance logic. It comprises:
\begin{itemize}
\item A machine learning risk scoring engine that computes a numerical risk 
score for each submitted transaction based on engineered features derived 
from address-level, transaction-level, and graph-level attributes.
\item A graph analytics module that constructs and analyses transaction 
graphs using a graph database, computing graph-level features including 
degree centrality, betweenness centrality, clustering coefficient, and 
community membership.
\item A sanctions screening module that checks transaction counterparties 
against maintained sanctions lists including OFAC (Office of Foreign Assets 
Control) and EU sanctions databases.
\item A policy engine that applies a deterministic compliance decision matrix 
to map computed risk scores to policy actions.
\item A state-adaptive transaction flow control module configured to compute 
one or more network instability metrics and to select a predetermined 
enforcement profile that determines one or more enforcement parameters used 
by the policy engine.
\end{itemize}

\textbf{Layer III --- Offline Training Layer.} This layer implements the 
machine learning model training pipeline. It comprises:
\begin{itemize}
\item A composite teacher ensemble for generating pseudo-labels on 
unlabelled data.
\item A student model training pipeline using the pseudo-labels.
\item Cross-validation and calibration procedures.
\end{itemize}

\textbf{Layer IV --- Infrastructure Layer.} This layer provides the 
persistent data storage, blockchain smart contracts, and cryptographic proof 
systems. It comprises:
\begin{itemize}
\item Blockchain smart contracts deployed on an Ethereum-compatible network 
that enforce compliance decisions deterministically.
\item A zero-knowledge proof subsystem (zkNAF).
\item Database systems for transaction records, graph data, and audit trails.
\item Distributed file storage for immutable audit records.
\end{itemize}

\subsection{Composite Teacher Weak Supervision Pipeline}

The offline training layer implements a novel weak supervision approach to 
address the scarcity of labelled fraud data in blockchain transactions. The 
composite teacher ensemble generates pseudo-labels as follows.

Three independent teacher models are constructed:
\begin{enumerate}[label=(\alph*)]
\item An autoencoder--gradient boosting teacher ($s_{\text{AE-XGB}}$) that 
trains an autoencoder on transaction features to learn a compressed 
representation, then trains a gradient boosting model on the reconstruction 
errors and latent features to estimate anomaly scores.

\item A rule-based FATF teacher ($s_{\text{FATF}}$) that encodes Financial 
Action Task Force red-flag indicators as deterministic scoring rules, 
including rapid fund movement patterns, round-number transactions, 
structuring behaviour, and high-risk jurisdiction involvement.

\item A graph-based teacher ($s_{\text{Graph}}$) that computes suspicion 
scores based on graph-structural properties of the transaction network, 
including proximity to known high-risk addresses, unusual connectivity 
patterns, and temporal clustering of transactions.
\end{enumerate}

The composite teacher score for each transaction $x$ is computed as:
\begin{equation}
s_{\text{teacher}}(x) = w_1 \cdot s_{\text{AE-XGB}}(x) + w_2 \cdot 
s_{\text{FATF}}(x) + w_3 \cdot s_{\text{Graph}}(x)
\end{equation}
where in the preferred embodiment, $w_1 = 0.4$, $w_2 = 0.3$, $w_3 = 0.3$.

The composite teacher scores are thresholded to generate binary 
pseudo-labels, which are then used to train a student model. In the 
preferred embodiment, the student model is an XGBoost gradient boosting 
classifier trained on 167 engineered features.

\subsection{Risk Score Calibration}

The raw output of the student model is recalibrated using a sigmoid 
recalibration function:
\begin{equation}
\hat{p}_{\text{cal}}(x) = \sigma\left(k \cdot (s(x) - P_{75})\right)
\end{equation}
where $\sigma$ is the logistic sigmoid function, $k$ is a steepness 
parameter (in the preferred embodiment, $k = 30$), and $P_{75}$ is the 75th 
percentile of the training score distribution. This recalibration compresses 
mid-range scores and amplifies boundary decisions, improving discrimination 
at the critical decision thresholds.

The calibrated probability is then mapped to an integer risk score on a 
0--1000 scale for use by the compliance decision matrix.

\subsection{Compliance Decision Matrix}

The policy engine applies a deterministic decision matrix that maps 
integer risk scores to policy actions:

\begin{center}
\begin{tabular}{lll}
\toprule
\textbf{Score Range} & \textbf{Policy Action} & \textbf{Smart Contract Effect} \\
\midrule
$[0, 200)$ & APPROVE & Transaction executed immediately \\
$[200, 400)$ & APPROVE with monitoring & Transaction executed; monitoring flag set \\
$[400, 600)$ & REVIEW & Transaction queued for manual review \\
$[600, 800)$ & ESCROW & Funds locked in escrow smart contract \\
$[800, 1000]$ & BLOCK & Transaction rejected on-chain \\
\bottomrule
\end{tabular}
\end{center}

These thresholds are fixed and public, ensuring that compliance decisions are 
reproducible, auditable, and deterministic. Every transaction receives 
identical treatment for identical risk scores, regardless of the time, 
venue, or identity of the transacting parties.

\subsection{State-Adaptive Transaction Flow Control (Adaptive Friction)}

In an enhanced embodiment, the compliance orchestration layer comprises a 
state-adaptive transaction flow control module (also referred to herein as an 
adaptive friction controller). The purpose of this module is to modulate the 
enforcement strictness of the deterministic compliance lifecycle in response 
to system-wide instability, while preserving determinism and auditability.

The module computes one or more network instability metrics from:
\begin{enumerate}[label=(\alph*)]
\item a transaction graph constructed from recent on-chain activity; and/or
\item aggregate transaction-risk dynamics comprising one or more of: 
risk score distribution shift, high-risk transaction rate, graph centrality 
concentration, and temporal burst indicators.
\end{enumerate}

Based on the computed instability metrics, the module selects an enforcement 
profile from a predetermined set of profiles (for example: \textit{stable}, 
\textit{stressed}, and \textit{critical}). Each profile deterministically 
defines one or more enforcement parameters used by the policy engine and/or 
smart contracts, including:
\begin{enumerate}[label=(\alph*)]
\item one or more risk score thresholds used by the decision matrix;
\item a transaction rate limit or queue delay for review actions;
\item an escrow duration or release condition window; and
\item additional proof requirements (e.g., requiring a risk-range proof or 
sanctions proof for specific classes of transactions).
\end{enumerate}

Accordingly, for a fixed system state (instability metrics and profile 
selection inputs), the selected enforcement profile and all resulting 
enforcement decisions are deterministic and reproducible.

\subsection{Zero-Knowledge Proof Subsystem (zkNAF)}

The zero-knowledge non-disclosing anti-fraud (zkNAF) subsystem enables 
participants to prove compliance properties without revealing underlying 
sensitive data. The subsystem implements three distinct proof circuits using 
the Groth16 zero-knowledge succinct non-interactive argument of knowledge 
(zk-SNARK) protocol:

\subsubsection{Sanctions Non-Membership Proof}

This circuit proves that a blockchain address is not a member of a set of 
sanctioned addresses, without revealing the address being verified. The 
circuit operates as follows:
\begin{enumerate}
\item The sanctioned address set is organised as a Merkle tree using the 
Poseidon hash function (chosen for its efficiency within arithmetic circuits).
\item The prover demonstrates possession of a valid Merkle non-membership 
proof---i.e., that the address does not appear at any leaf of the tree.
\item The verifier checks the proof against the published Merkle root of the 
current sanctions list without learning the prover's address.
\end{enumerate}

\subsubsection{Risk Range Proof}

This circuit proves that a computed risk score falls within a specified 
acceptable range, without revealing the exact score. Formally, the prover 
demonstrates:
\begin{equation}
\text{score}_{\min} \leq s(x) \leq \text{score}_{\max}
\end{equation}
where $s(x)$ is the computed risk score and $[\text{score}_{\min}, 
\text{score}_{\max}]$ is the acceptable range, without revealing $s(x)$.

\subsubsection{KYC Verification Proof}

This circuit proves that a participant has completed Know Your Customer 
verification with an accredited provider without revealing the participant's 
identity data. The circuit verifies a digital commitment to the KYC status 
and provider attestation.

\subsubsection{Proof Lifecycle}

Each zero-knowledge proof has a time-to-live of 24 hours, after which it 
must be regenerated. Proofs are verified on-chain by a dedicated verifier 
smart contract that stores verification results. The verification function 
accepts the proof components ($\pi_a$, $\pi_b$, $\pi_c$) and public inputs, 
performs the bilinear pairing check, and records the result with a timestamp.

\subsection{Smart Contract Architecture}

The infrastructure layer comprises a plurality of smart contracts deployed 
on an Ethereum-compatible blockchain network. The principal contracts 
include:

\begin{itemize}
\item \textbf{Core Contract:} Receives transaction risk scores from the 
compliance orchestration layer and registers compliance decisions on-chain. 
Implements the deterministic decision matrix as an on-chain mapping from 
score ranges to policy actions.

\item \textbf{Escrow Contract:} Implements a time-locked escrow mechanism 
for transactions assigned the ESCROW policy action. Funds are held for a 
configurable review period during which authorised reviewers may release the 
funds (if the review outcome is favourable) or initiate a dispute resolution 
process.

\item \textbf{ZK Verifier Contract:} Verifies Groth16 zero-knowledge proofs 
on-chain. Accepts proof components and public inputs, performs the 
elliptic-curve pairing verification, and stores the binary verification 
result with a timestamp for subsequent reference by other contracts.

\item \textbf{Cross-Chain Bridge Contract:} Extends compliance enforcement 
to cross-chain transactions by implementing a bridge mechanism that requires 
compliance verification on both the source and destination chains before 
releasing bridged funds.

\item \textbf{Dispute Resolution Contract:} Implements a structured process 
for resolving disputes arising from escrowed transactions, including 
escalation procedures and release conditions.
\end{itemize}

\subsection{Feature Engineering}

The machine learning risk scoring engine operates on 167 engineered features 
derived from three categories:

\begin{enumerate}[label=(\alph*)]
\item \textbf{Address-level features:} Transaction count, total value 
transferred, mean and variance of transaction amounts, age of address, 
number of unique counterparties, and temporal activity patterns.

\item \textbf{Transaction-level features:} Transaction value, gas 
parameters, input data characteristics, number of internal transactions, 
and token transfer indicators.

\item \textbf{Graph-level features:} Degree centrality, betweenness 
centrality, closeness centrality, clustering coefficient, PageRank score, 
community membership identifier, and local graph density.
\end{enumerate}

\subsection{Performance Characteristics}

In the preferred embodiment, trained on 2.64 million Ethereum transactions, 
the system achieves the following performance:

\begin{center}
\begin{tabular}{ll}
\toprule
\textbf{Metric} & \textbf{Value} \\
\midrule
Meta-Ensemble ROC-AUC (test) & 1.0000 \\
Meta-Ensemble ROC-AUC (5-fold CV) & $0.9965 \pm 0.0049$ \\
Meta-Ensemble PR-AUC & 0.9999 \\
Transaction-Level XGBoost ROC-AUC & 0.9878 \\
Transaction-Level XGBoost F1 & 0.901 \\
Cross-domain mean ROC-AUC (9 datasets) & 0.9962 \\
Escrow smart contract gas cost & 287,890 gas \\
\bottomrule
\end{tabular}
\end{center}

% ============================================================
\section{CLAIMS}
% ============================================================

\noindent\textit{The following claims define the scope of protection
sought. Where a claim refers to another claim, the dependency is by claim
number. Features of any dependent claim may be combined with any
independent claim or with other dependent claims.}
\vspace{0.8em}

\begin{enumerate}[label=\textbf{\arabic*.}]

% --- Independent System Claim ---
\item A computer-implemented transaction compliance system for decentralised 
finance networks, the system comprising:
\begin{enumerate}[label=(\alph*)]
\item an interface layer configured to receive transaction requests from 
client applications via software development kits or application programming 
interfaces;

\item a compliance orchestration layer comprising:
\begin{enumerate}[label=(\roman*)]
\item a machine learning risk scoring engine configured to compute a 
numerical risk score for each received transaction request based on a 
plurality of engineered features derived from address-level, 
transaction-level, and graph-level attributes;
\item a graph analytics module configured to construct and analyse 
transaction graphs and compute graph-level features;
\item a sanctions screening module configured to check transaction 
counterparties against maintained sanctions databases; and
\item a state-adaptive transaction flow control module configured to compute 
one or more network instability metrics from recent transaction activity and 
to select, from a predetermined set of enforcement profiles, an enforcement 
configuration comprising:
\begin{enumerate}[label=(\alph*)]
\item a risk threshold profile defining numerical thresholds for mapping risk 
scores to policy actions; and
\item a proof-requirement policy defining, for at least one risk score tier 
or transaction class, one or more cryptographic proofs that are required to 
be verified as a prerequisite to settlement;
\end{enumerate}
and
\item a deterministic policy engine configured to:
\begin{enumerate}[label=(\alph*)]
\item map computed risk scores to policy actions selected from the group 
consisting of: approve, approve with monitoring, review, escrow, and block, 
according to numerical thresholds defined by the selected risk threshold 
profile; and
\item enforce the selected proof-requirement policy by causing verification 
of the required cryptographic proofs on-chain prior to settlement, and by 
assigning a non-approval policy action when the required cryptographic proofs 
are missing or fail verification;
\end{enumerate}
\end{enumerate}

\item an offline training layer comprising a composite teacher ensemble 
configured to generate pseudo-labels from unlabelled blockchain transaction 
data using a weighted combination of at least two independent teacher 
models, and a student model trained on the generated pseudo-labels;

\item an infrastructure layer comprising a plurality of smart contracts 
deployed on a blockchain network, the smart contracts being configured to 
enforce the policy actions determined by the policy engine, including:
\begin{enumerate}[label=(\roman*)]
\item a core contract configured to register compliance decisions on the 
blockchain;
\item an escrow contract configured to hold funds in time-locked escrow for 
transactions assigned the escrow policy action; and
\item a zero-knowledge proof verifier contract configured to verify 
cryptographic proofs on-chain;
\end{enumerate}
\end{enumerate}
wherein the system is configured to enforce compliance decisions 
\emph{before} the corresponding transaction is settled on the blockchain 
network.

% --- Independent Method Claim ---
\item A computer-implemented method of enforcing compliance in decentralised 
finance transactions, the method comprising the steps of:
\begin{enumerate}[label=(\alph*)]
\item receiving a transaction request comprising at least a sender address, 
a recipient address, and a transaction value;
\item computing a risk score for the transaction request by applying a 
trained machine learning model to a feature vector comprising address-level 
features, transaction-level features, and graph-level features derived from 
the transaction request and associated blockchain data;
\item calibrating the computed risk score using a sigmoid recalibration 
function anchored at a percentile of a training score distribution;
\item computing one or more network instability metrics from recent 
transaction activity and selecting an enforcement profile from a 
predetermined set of enforcement profiles, the selected enforcement profile 
deterministically defining:
\begin{enumerate}[label=(\roman*)]
\item a risk threshold profile defining numerical thresholds for a 
deterministic decision matrix; and
\item a proof-requirement policy defining one or more cryptographic proofs 
to be verified on-chain as a prerequisite to settlement for at least one 
risk score tier or transaction class;
\end{enumerate}
\item mapping the calibrated risk score to a policy action selected from a 
predetermined set of policy actions according to the deterministic decision 
matrix using numerical thresholds defined by the selected risk threshold 
profile;
\item enforcing the proof-requirement policy by verifying the one or more 
cryptographic proofs on-chain prior to settlement, and assigning a non-
approval policy action when the required cryptographic proofs are missing or 
fail verification;
\item enforcing the policy action by executing a corresponding smart 
contract function on a blockchain network, prior to settlement of the 
transaction; and
\item recording the compliance decision and policy action on the blockchain 
as an immutable audit record.
\end{enumerate}

% --- Dependent Claims ---
\item A system according to claim 1, wherein the composite teacher ensemble 
comprises:
\begin{enumerate}[label=(\alph*)]
\item an autoencoder--gradient boosting teacher configured to compute 
anomaly scores based on reconstruction errors of an autoencoder applied to 
transaction features;
\item a rule-based teacher configured to compute suspicion scores based on 
Financial Action Task Force (FATF) red-flag indicators; and
\item a graph-based teacher configured to compute suspicion scores based on 
graph-structural properties of the transaction network;
\end{enumerate}
and wherein the composite teacher score is a weighted sum of the outputs of 
the at least three teacher models.

\item A system according to claim 1, further comprising a zero-knowledge 
proof subsystem configured to generate and verify one or more of:
\begin{enumerate}[label=(\alph*)]
\item a sanctions non-membership proof demonstrating that an address is not 
a member of a set of sanctioned addresses represented as a Merkle tree, 
without revealing the address;
\item a risk range proof demonstrating that a risk score falls within a 
specified range, without revealing the exact risk score; and
\item a KYC verification proof demonstrating that a participant has 
completed identity verification, without revealing identity data.
\end{enumerate}

\item A system according to claim 4, wherein the Merkle tree of sanctioned 
addresses is constructed using the Poseidon hash function, and wherein the 
zero-knowledge proofs are Groth16 zk-SNARK proofs with a time-to-live after 
which the proofs must be regenerated.

\item A method according to claim 2, wherein the step of computing a risk 
score comprises:
\begin{enumerate}[label=(\alph*)]
\item computing address-level features comprising transaction count, total 
value transferred, mean and variance of transaction amounts, address age, 
number of unique counterparties, and temporal activity patterns;
\item computing transaction-level features comprising transaction value, gas 
parameters, input data characteristics, and token transfer indicators;
\item computing graph-level features by constructing a transaction graph and 
computing degree centrality, betweenness centrality, clustering coefficient, 
PageRank score, and community membership; and
\item applying the trained machine learning model to the combined feature 
vector.
\end{enumerate}

\item A system according to claim 1, wherein the deterministic decision 
matrix maps risk scores as follows:
\begin{enumerate}[label=(\alph*)]
\item scores below a first threshold to an approve action;
\item scores between the first threshold and a second threshold to an 
approve-with-monitoring action;
\item scores between the second threshold and a third threshold to a review 
action;
\item scores between the third threshold and a fourth threshold to an escrow 
action; and
\item scores at or above the fourth threshold to a block action.
\end{enumerate}

\item A method according to claim 2, wherein verifying the one or more 
cryptographic proofs comprises verifying one or more zero-knowledge proofs 
submitted with the transaction request on-chain, wherein successful 
verification is a prerequisite for approving the transaction request.

\item A system or method according to claim 1 or claim 2, wherein the one or 
more network instability metrics comprise at least one of: a spectral radius 
or centrality concentration of a transaction graph computed over a sliding 
time window, a temporal burst indicator computed from a transaction count 
process, or a shift of a calibrated risk score distribution relative to a 
reference distribution.

\item A system or method according to claim 1 or claim 2, wherein selecting 
the enforcement profile comprises selecting one profile from a discrete set 
of profiles, and wherein each profile deterministically defines a distinct 
set of risk score thresholds for the deterministic decision matrix.

\item A system or method according to claim 1 or claim 2, wherein the proof-
requirement policy comprises, for at least one of the enforcement profiles, 
requiring a sanctions non-membership proof, a risk range proof, a KYC 
verification proof, or a combination thereof.

\item A system or method according to claim 1 or claim 2, wherein selecting 
the enforcement profile further deterministically defines at least one of: a 
transaction rate limit, a queue delay, an escrow duration, an escrow release 
condition window, or a maximum transaction value limit.

\item A system according to claim 1, further comprising a cross-chain bridge 
contract configured to extend compliance enforcement to cross-chain 
transactions by requiring compliance verification on both a source chain and 
a destination chain before releasing bridged funds.

\item A method according to claim 2, wherein the trained machine learning 
model is a student model trained using pseudo-labels generated by a 
composite teacher ensemble operating on unlabelled blockchain transaction 
data, the composite teacher ensemble comprising at least an unsupervised 
anomaly detection model, a rule-based model encoding regulatory red-flag 
indicators, and a graph-structural model.

\item A system or method according to any preceding claim, wherein the 
sigmoid recalibration function is defined as:
\[
\hat{p}_{\text{cal}}(x) = \sigma\left(k \cdot (s(x) - P_{75})\right)
\]
where $\sigma$ is the logistic sigmoid, $k$ is a steepness parameter, $s(x)$ 
is the raw model output, and $P_{75}$ is the 75th percentile of the training 
score distribution.

\end{enumerate}

% ============================================================
\section{DRAWINGS}
% ============================================================

\textit{[To be filed separately: System architecture diagram showing the 
four layers; flow diagram of the compliance decision process; schematic of 
the zkNAF proof generation and verification pipeline; block diagram of the 
composite teacher weak supervision pipeline.]}

% ============================================================
\newpage
\section*{ABSTRACT}
% ============================================================

A computer-implemented system enforces regulatory compliance in decentralised
finance transactions before settlement. An interface layer receives
transaction requests. A compliance orchestration layer comprising a machine
learning risk scoring engine, graph analytics module, sanctions screening
module and state-adaptive transaction flow control module computes network
instability metrics and selects an enforcement profile; the selected profile
deterministically defines a risk threshold profile for the compliance decision
matrix and a proof-requirement policy specifying cryptographic proofs to be
verified on-chain before settlement. A deterministic policy engine maps
calibrated risk scores to approve, review, escrow or block actions using
thresholds defined by the selected profile. Blockchain-deployed smart
contracts enforce those actions before transaction settlement, with an
immutable audit trail. A zero-knowledge proof subsystem enables
privacy-preserving verification of sanctions non-membership, risk score range
and KYC status without revealing underlying data.

\end{document}
